% Propietario conservado: /Users/ruben/Documents/New project/output/EXTREMA_MEDIA_RAZON_RECIPROCIDAD_ES_EN_20260918/ARTICULO/ES/source/sections/realizacion_lc.tex
\subsection{Realización inductiva--capacitiva de la forma de acción}
\label{x_reciprocidad:nuclear:realizacion-lc}

La forma de acción ya construida admite una realización eléctrica en
una carta provista de un reloj \(\omega>0\) y una impedancia de
referencia \(Z_*>0\). Estas dos magnitudes conservan sus unidades y
acompañan al parámetro de forma \(\eta\). Se definen
\begin{equation}
 L_\eta=\frac{Z_*}{\omega}e^{-\eta},\qquad
 C_\eta=\frac{e^\eta}{Z_*\omega},\qquad
 Z_\eta=\sqrt{L_\eta/C_\eta}.
 \label{x_reciprocidad:hmtII:eq:ii-elipse-lc}
\end{equation}
Aquí \(C_\eta\) denota capacitancia, distinta del contraángulo \(C^*\).

\begin{proposition}[Conservación del producto dinámico]
\label{x_reciprocidad:nuclear:lc-producto}
La realización \eqref{x_reciprocidad:hmtII:eq:ii-elipse-lc} satisface
\begin{equation}
 L_\eta C_\eta=\omega^{-2},\qquad
 \frac{Z_\eta}{Z_*}=e^{-\eta}=\frac{b_S}{a_S},
 \label{x_reciprocidad:hmtII:eq:ii-elipse-impedancia}
\end{equation}
donde \(a_S=\sqrt{2\hbar}\,e^{\eta/2}\) y
\(b_S=\sqrt{2\hbar}\,e^{-\eta/2}\) son los semiejes de
\cref{x_reciprocidad:iv:ellipse:semiejes}.
Con \(Q=\sqrt{Z_*}q\) y \(P=\Phi/\sqrt{Z_*}\), la energía toma la forma
\[
 H_{LC}=\frac{q^2}{2C_\eta}+\frac{\Phi^2}{2L_\eta}
       =\frac\omega2(e^{-\eta}Q^2+e^\eta P^2).
\]
Sobre \((Q,P)=(a_S\cos\theta,b_S\sin\theta)\), las contribuciones
eléctrica y magnética son \(\hbar\omega\cos^2\theta\) y
\(\hbar\omega\sin^2\theta\); la energía total es \(\hbar\omega\).
\end{proposition}
\begin{proof}
La multiplicación de \(L_\eta\) y \(C_\eta\) cancela los factores
\(Z_*\) y \(e^{\pm\eta}\), mientras su cociente es
\(Z_*^2e^{-2\eta}\). La raíz positiva prueba la identidad de impedancia.
Sustituir \(q=Q/\sqrt{Z_*}\) y \(\Phi=\sqrt{Z_*}P\) da la forma
de \(H_{LC}\). Los cuadrados de los semiejes cancelan los factores
de forma restantes; la suma utiliza \(\cos^2\theta+\sin^2\theta=1\).
\end{proof}

La convención hamiltoniana
\(\dot Q=\partial_PH\), \(\dot P=-\partial_QH\) asigna
\(\dot\theta=-\omega\) a esta parametrización. La forma, el reloj
y la carta dimensional permanecen diferenciados. La respuesta
constitutiva del vacío y la impedancia \(Z_\eta\) quedan relacionadas
por la inversión de canales de \cref{x_reciprocidad:iii:prop:forma-LC}; cada una
conserva su definición y su normalización.
